% claim.tex -- AI4Math claim of record (AI-generated, 2026-09-29; instantiated
% from harness/templates/claim/claim.tex per SOP 08b).
% Machine-readable theorem anchors are the "% LEAN:" lines; scripts/paper-lint.sh
% and the claim audit check them. All Lean listings are verbatim, byte-identical
% contiguous excerpts of frozen artifacts of tasks/20260929-l14c-graphside-pipeline:
%   statements : formalized/statement.lean (124 lines, sha256 d186ff16...)
%   proof file : attempts/phase0/r2/split-identity.lean (429 lines, sha256 d2efac51...)
% Line ranges used: statement 45-62 (base defs); final 215-259 (general-lemma
% head), 1-429 (whole file). The frozen statement file's three literal "sorry"
% occurrences (two bridge-theorem proof placeholders and one head-comment
% mention) all lie outside the excerpts, as required by the paper-lint listing
% gate. Verification: audit/listing-verify.txt (audit/inspect-tex.py --verify).
\documentclass[11pt]{article}

\usepackage[T1]{fontenc}
\usepackage[utf8]{inputenc}
\usepackage{amsmath,amssymb,amsthm}
\usepackage{graphicx}
\usepackage{hyperref}
\usepackage{xurl}
\setlength{\emergencystretch}{3em}
% lstlean.tex — Lean style for the listings package.
% Lineage: Jeremy Avigad (2015), adapted from Assia Mahboubi's SSR style;
% inspired by Olivier Verdier's unixode.sty for the Unicode literate table.
% No CTAN package or upstream repo exists; papers carry this file in their
% source package (cap set arXiv:1907.01449, sphere eversion arXiv:2210.07746).
% AI4Math adaptation (AI-generated, 2026-09-26): sorry/admit/sorryAx elevated
% to a dedicated error tier, matching the pipeline's 0-sorry constitution.
\usepackage{listings}
\usepackage{xcolor}

\definecolor{keywordcolor}{rgb}{0.7, 0.1, 0.1}   % red keywords
\definecolor{tacticcolor}{rgb}{0.0, 0.1, 0.6}    % blue tactics
\definecolor{commentcolor}{rgb}{0.4, 0.4, 0.4}   % grey comments
\definecolor{symbolcolor}{rgb}{0.0, 0.1, 0.6}    % blue symbols
\definecolor{sortcolor}{rgb}{0.1, 0.5, 0.1}      % green sorts
\definecolor{errorcolor}{rgb}{0.6, 0, 0}         % dark red: sorry/admit
\definecolor{stringcolor}{rgb}{0.5, 0.1, 0.5}    % purple strings

\lstdefinelanguage{lean}{
  % Anything between $ becomes LaTeX math mode
  mathescape=false,
  % Comments may or not include Latex commands
  texcl=false,
  % keywords, list taken from lean-syntax.txt
  morekeywords=[1]{import, prelude, protected, private, noncomputable,
    definition, meta, renaming, hiding, parameter, parameters, begin,
    constant, constants, lemma, variable, variables, theory, print,
    theorem, example, open, section, namespace, universe, universes, end,
    modifier, modifiers, using, export, exposing, axiom, inductive,
    coinductive, with, structure, class, instance, attribute, local, where,
    by, have, show, let, in, if, then, else, match, do, for, fun, assume,
    from, take, calc, include, omit, infix, infixl, infixr, notation,
    macro, syntax, elab, set_option, initialize, deriving, opaque, abbrev,
    mutual, partial, scoped, termination_by, decreasing_by},
  % Sorts
  morekeywords=[2]{Sort, Type, Prop},
  % tactics, list taken from lean-syntax.txt (tactic keywords)
  morekeywords=[3]{intro, intros, apply, exact, refine, rfl, simp, simp_all,
    simpa, simp_rw, rw, rewrite, erewrite, ring, ring_nf, omega, decide,
    norm_num, norm_cast, induction, cases, rcases, obtain, constructor,
    ext, funext, conv, congr, field_simp, linarith, nlinarith, positivity,
    tauto, trivial, aesop, use, by_cases, by_contra, push_neg, contrapose,
    symm, trans, assumption, contradiction, exfalso, specialize,
    generalize, revert, clear, rename_i, subst, unfold, delta, change,
    convert, split_ifs, interval_cases, fin_cases, zify, gcongr,
    nth_rewrite, suffices, generalizing, at, only, native_decide},
  % warnings - constitution tier: must never survive into a final proof
  morekeywords=[4]{sorry, admit, sorryAx},
  literate=
    {α}{{\ensuremath{\alpha}}}1
    {β}{{\ensuremath{\beta}}}1
    {γ}{{\ensuremath{\gamma}}}1
    {δ}{{\ensuremath{\delta}}}1
    {ε}{{\ensuremath{\varepsilon}}}1
    {ζ}{{\ensuremath{\zeta}}}1
    {η}{{\ensuremath{\eta}}}1
    {θ}{{\ensuremath{\theta}}}1
    {ι}{{\ensuremath{\iota}}}1
    {κ}{{\ensuremath{\kappa}}}1
    {λ}{{\ensuremath{\lambda}}}1
    {μ}{{\ensuremath{\mu}}}1
    {ν}{{\ensuremath{\nu}}}1
    {ξ}{{\ensuremath{\xi}}}1
    {π}{{\ensuremath{\pi}}}1
    {ρ}{{\ensuremath{\rho}}}1
    {σ}{{\ensuremath{\sigma}}}1
    {τ}{{\ensuremath{\tau}}}1
    {υ}{{\ensuremath{\upsilon}}}1
    {φ}{{\ensuremath{\varphi}}}1
    {χ}{{\ensuremath{\chi}}}1
    {ψ}{{\ensuremath{\psi}}}1
    {ω}{{\ensuremath{\omega}}}1
    {Γ}{{\ensuremath{\Gamma}}}1
    {Δ}{{\ensuremath{\Delta}}}1
    {Θ}{{\ensuremath{\Theta}}}1
    {Λ}{{\ensuremath{\Lambda}}}1
    {Ξ}{{\ensuremath{\Xi}}}1
    {Π}{{\ensuremath{\Pi}}}1
    {Σ}{{\ensuremath{\Sigma}}}1
    {Φ}{{\ensuremath{\Phi}}}1
    {Ψ}{{\ensuremath{\Psi}}}1
    {Ω}{{\ensuremath{\Omega}}}1
    {ℵ}{{\ensuremath{\aleph}}}1
    {∀}{{\ensuremath{\forall}}}1
    {∃!}{{\ensuremath{\exists}!}}1
    {∃}{{\ensuremath{\exists}}}1
    {∈}{{\ensuremath{\in}}}1
    {∉}{{\ensuremath{\notin}}}1
    {∘}{{\ensuremath{\circ}}}1
    {∩}{{\ensuremath{\cap}}}1
    {∪}{{\ensuremath{\cup}}}1
    {⊆}{{\ensuremath{\subseteq}}}1
    {⊂}{{\ensuremath{\subset}}}1
    {⊇}{{\ensuremath{\supseteq}}}1
    {⊃}{{\ensuremath{\supset}}}1
    {⊄}{{\ensuremath{\nsubseteq}}}1
    {∅}{{\ensuremath{\emptyset}}}1
    {∧}{{\ensuremath{\wedge}}}1
    {∨}{{\ensuremath{\vee}}}1
    {¬}{{\ensuremath{\neg}}}1
    {⊤}{{\ensuremath{\top}}}1
    {⊥}{{\ensuremath{\bot}}}1
    {↔}{{\ensuremath{\leftrightarrow}}}1
    {→}{{\ensuremath{\rightarrow}}}1
    {←}{{\ensuremath{\leftarrow}}}1
    {↑}{{\ensuremath{\uparrow}}}1
    {⇒}{{\ensuremath{\Rightarrow}}}1
    {⇐}{{\ensuremath{\Leftarrow}}}1
    {⟹}{{\ensuremath{\Longrightarrow}}}1
    {↦}{{\ensuremath{\mapsto}}}1
    {≤}{{\ensuremath{\leq}}}1
    {≥}{{\ensuremath{\geq}}}1
    {≠}{{\ensuremath{\neq}}}1
    {≈}{{\ensuremath{\approx}}}1
    {≡}{{\ensuremath{\equiv}}}1
    {≃}{{\ensuremath{\simeq}}}1
    {≅}{{\ensuremath{\cong}}}1
    {×}{{\ensuremath{\times}}}1
    {·}{{\ensuremath{\cdot}}}1
    {±}{{\ensuremath{\pm}}}1
    {⊕}{{\ensuremath{\oplus}}}1
    {⊗}{{\ensuremath{\otimes}}}1
    {⋆}{{\ensuremath{\star}}}1
    {▸}{{\ensuremath{\blacktriangleright}}}1
    {⟨}{{\ensuremath{\langle}}}1
    {⟩}{{\ensuremath{\rangle}}}1
    {⦃}{{\ensuremath{\{\!\mid}}}1
    {⦄}{{\ensuremath{\mid\!\}}}}1
    {⟦}{{\ensuremath{[\![}}}1
    {⟧}{{\ensuremath{]\!]}}}1
    {⌊}{{\ensuremath{\lfloor}}}1
    {⌋}{{\ensuremath{\rfloor}}}1
    {⌈}{{\ensuremath{\lceil}}}1
    {⌉}{{\ensuremath{\rceil}}}1
    {√}{{\ensuremath{\surd}}}1
    {∣}{{\ensuremath{\mid}}}1
    {∤}{{\ensuremath{\nmid}}}1
    {∎}{{\ensuremath{\blacksquare}}}1
    {⋯}{{\ensuremath{\cdots}}}1
    {…}{{\ensuremath{\ldots}}}1
    {∑}{{\ensuremath{\sum}}}1
    {∏}{{\ensuremath{\prod}}}1
    {⁻¹}{{\ensuremath{^{-1}}}}1
    {¹}{{\ensuremath{^1}}}1
    {²}{{\ensuremath{^2}}}1
    {³}{{\ensuremath{^3}}}1
    {ⁿ}{{\ensuremath{^n}}}1
    {ᵐ}{{\ensuremath{^m}}}1
    {₀}{{\ensuremath{_0}}}1
    {₁}{{\ensuremath{_1}}}1
    {₂}{{\ensuremath{_2}}}1
    {₃}{{\ensuremath{_3}}}1
    {₄}{{\ensuremath{_4}}}1
    {₅}{{\ensuremath{_5}}}1
    {ₙ}{{\ensuremath{_n}}}1
    {ₘ}{{\ensuremath{_m}}}1
    {ₖ}{{\ensuremath{_k}}}1
    {ᵢ}{{\ensuremath{_i}}}1
    {ⱼ}{{\ensuremath{_j}}}1
    {ℤ}{{\ensuremath{\mathbb{Z}}}}1
    {ℕ}{{\ensuremath{\mathbb{N}}}}1
    {ℚ}{{\ensuremath{\mathbb{Q}}}}1
    {ℝ}{{\ensuremath{\mathbb{R}}}}1
    {ℂ}{{\ensuremath{\mathbb{C}}}}1
    {𝔽}{{\ensuremath{\mathbb{F}}}}1
    {𝒫}{{\ensuremath{\mathcal{P}}}}1
    {∗}{{\ensuremath{\ast}}}1
    {•}{{\ensuremath{\bullet}}}1,
  % Comments
  morecomment=[l]{--},
  morecomment=[s]{/-}{-/},
  morestring=[b]",
  stringstyle=\color{stringcolor},
  keepspaces=true,
  keywordstyle=[1]\color{keywordcolor}\bfseries,
  keywordstyle=[2]\color{sortcolor}\bfseries,
  keywordstyle=[3]\color{tacticcolor}\bfseries,
  keywordstyle=[4]\color{errorcolor}\bfseries\underline,
  escapeinside={(*@}{@*)},
  columns=fullflexible,
  basicstyle=\ttfamily\small,
  commentstyle=\color{commentcolor}\itshape,
  breaklines=true,
  showstringspaces=false,
  frame=single,
  framesep=2mm,
  xleftmargin=4mm,
  numbers=left,
  numberstyle=\tiny\color{commentcolor},
  numbersep=6pt,
}

% Inline Lean code in prose: \lean{Int.ModEq.pow}
\newcommand{\lean}[1]{\lstinline[language=lean,basicstyle=\ttfamily\normalsize]|#1|}

\lstset{language=lean}

% --- local rendering patch (2026-09-29) ---
% tectonic/XeTeX does not apply the listings literate table to multi-byte
% characters, so the Unicode set used by the listings of THIS note is mapped to
% math glyphs as active characters via the standard \lccode/\lowercase idiom
% (ASCII-only source, no extra packages). Restricted to symbols outside xeCJK's
% own punctuation table: 00B7 (cdot) and 2014 (em-dash) also occur in the frozen
% sources but are special punctuation that xeCJK sets up at load time -- making
% them active first aborts the run with "Control sequence ... already defined"
% (observed 2026-09-27 in the pendant note); they route through xeCJK + SimSun
% instead, as do all CJK ideographs and fullwidth/CJK punctuation (FF0C, FF08,
% FF09, FF1A, FF1B, 3001, 3002, 300C, 300D).
% New relative to the pendant note (they occur in this note's sources): 03B1/
% 03B2 (Greek letters, missing shapes in lmmono), 2115/2124 (N/Z blackboard),
% 2208 (in), 2209 (notin), 2260 (ne), 2081/2082 (subscript digits), 2191
% (uparrow), 2194 (leftrightarrow), 21A5 (upwards from bar), 25A1 (box), 22A2
% (turnstile).
% Rendering only: listing source bytes are untouched, so the byte-exactness
% checks are unaffected, and the \ifdefined guard leaves the pdflatex path
% identical to the template.
\ifdefined\XeTeXversion
  {\lccode`\~="03B1 \lowercase{\global\catcode`~=\active \global\def~{\ensuremath{\alpha}}}}
  {\lccode`\~="03B2 \lowercase{\global\catcode`~=\active \global\def~{\ensuremath{\beta}}}}
  {\lccode`\~="2081 \lowercase{\global\catcode`~=\active \global\def~{\ensuremath{_{1}}}}}
  {\lccode`\~="2082 \lowercase{\global\catcode`~=\active \global\def~{\ensuremath{_{2}}}}}
  {\lccode`\~="2115 \lowercase{\global\catcode`~=\active \global\def~{\ensuremath{\mathbb{N}}}}}
  {\lccode`\~="2124 \lowercase{\global\catcode`~=\active \global\def~{\ensuremath{\mathbb{Z}}}}}
  {\lccode`\~="2190 \lowercase{\global\catcode`~=\active \global\def~{\ensuremath{\leftarrow}}}}
  {\lccode`\~="2191 \lowercase{\global\catcode`~=\active \global\def~{\ensuremath{\uparrow}}}}
  {\lccode`\~="2192 \lowercase{\global\catcode`~=\active \global\def~{\ensuremath{\rightarrow}}}}
  {\lccode`\~="2194 \lowercase{\global\catcode`~=\active \global\def~{\ensuremath{\leftrightarrow}}}}
  {\lccode`\~="21A5 \lowercase{\global\catcode`~=\active \global\def~{\ensuremath{\upharpoonleft}}}}
  {\lccode`\~="2200 \lowercase{\global\catcode`~=\active \global\def~{\ensuremath{\forall}}}}
  {\lccode`\~="2203 \lowercase{\global\catcode`~=\active \global\def~{\ensuremath{\exists}}}}
  {\lccode`\~="2208 \lowercase{\global\catcode`~=\active \global\def~{\ensuremath{\in}}}}
  {\lccode`\~="2209 \lowercase{\global\catcode`~=\active \global\def~{\ensuremath{\notin}}}}
  {\lccode`\~="2212 \lowercase{\global\catcode`~=\active \global\def~{\ensuremath{-}}}}
  {\lccode`\~="2227 \lowercase{\global\catcode`~=\active \global\def~{\ensuremath{\wedge}}}}
  {\lccode`\~="2228 \lowercase{\global\catcode`~=\active \global\def~{\ensuremath{\vee}}}}
  {\lccode`\~="2260 \lowercase{\global\catcode`~=\active \global\def~{\ensuremath{\neq}}}}
  {\lccode`\~="2264 \lowercase{\global\catcode`~=\active \global\def~{\ensuremath{\leq}}}}
  {\lccode`\~="2265 \lowercase{\global\catcode`~=\active \global\def~{\ensuremath{\geq}}}}
  {\lccode`\~="22A2 \lowercase{\global\catcode`~=\active \global\def~{\ensuremath{\vdash}}}}
  {\lccode`\~="25A1 \lowercase{\global\catcode`~=\active \global\def~{\ensuremath{\square}}}}
  {\lccode`\~="25B8 \lowercase{\global\catcode`~=\active \global\def~{\ensuremath{\blacktriangleright}}}}
  {\lccode`\~="27E8 \lowercase{\global\catcode`~=\active \global\def~{\ensuremath{\langle}}}}
  {\lccode`\~="27E9 \lowercase{\global\catcode`~=\active \global\def~{\ensuremath{\rangle}}}}
\fi
\ifdefined\XeTeXversion
  \usepackage{xeCJK}
  \setCJKmainfont{SimSun}
\fi

\newtheorem{theorem}{Theorem}
\theoremstyle{definition}
\newtheorem{lemma}{Lemma}

\title{An edge-deletion matching-split identity for finite simple graphs, and
its instantiation on the prism graph\\ (Lean 4 machine-checked claim of
record)}
\author{Aurora0134\thanks{Contact: via the GitHub account Aurora0134; ORCID
placeholder --- to be filled in by the author before any upload. No personal
information is carried in this note.}}
\date{September 29, 2026}

\begin{document}
\maketitle

\begin{abstract}
For any finite simple graph $G$ with a decidable adjacency relation and any
edge $ab$ of $G$, the number of matchings of $G$ splits as the number of
matchings of $G$ with the edge $ab$ deleted, plus the number of matchings of
the subgraph induced on the complement of the two endpoints: every matching
either avoids $ab$ or contains it, and a matching containing $ab$ uses no other
edge at either endpoint. This classical counting device of the matching
polynomial / Hosoya-index literature has, to the knowledge reachable by this
pipeline's novelty search, no Lean formalisation: mathlib (v4.34.0) has
matching predicates but no matching-count API at all, and no same-form
statement was found in the public Lean ecosystem. We report a complete,
machine-checked Lean 4 formalisation of the identity as a general lemma about
arbitrary graphs, together with its instantiation on the prism graph
$C_n\square P_2$ (delete one spoke, endpoints kept) where the induced factor is
a path $P_{n-1}$, giving $\mu(C_n\square P_2)=\mu(\text{prism minus one
spoke})+\lfloor \cdots \rfloor$ as the first structural step of the pipeline's
defective-prism program. Both theorems compile with no \texttt{sorry}, and
\texttt{\#print axioms} reports only \{propext, Quot.sound,
Classical.choice\}. The identity is asserted as a \emph{first Lean
formalisation of a known classical tool}, never as new mathematics. This note
is a priority claim of record, not a full paper.
\end{abstract}

\section{Introduction}

One lemma about arbitrary finite simple graphs and one instantiation on the
prism graph are reported here. The lemma is the standard edge-deletion split
of matching counts: for a graph $G$ and an edge $e=ab$,
\[
\mu(G)\;=\;\mu(G-e)\;+\;\mu(G-\{a,b\}),
\]
where $\mu$ counts all matchings (the empty one included) and $G-\{a,b\}$
denotes the graph induced on the complement of $\{a,b\}$. It follows from the
partition of matchings into those avoiding and those containing $e$; a
matching containing $e$ cannot use any other edge incident to $a$ or $b$, and
after removing $e$ the remainder is exactly a matching of the induced graph on
the remaining vertices. In the matching polynomial and Hosoya-index
literature this split is the classical per-edge device behind essentially
every computed matching count \cite{godsil-gutman}; the prism graph's own OEIS
entry A102080 \cite{a102080} is indexed under the recurrence it feeds. The
pipeline's novelty search (Section~\ref{sec:novelty}) found no Lean
formalisation of the identity anywhere in the public ecosystem, and mathlib
has no matching-count API to state it against \cite{mathlib}; the identity is
therefore claimed at the tier \emph{first formalisation (known unformalised)},
with no claim of new mathematics.

The proofs were produced by the AI4Math pipeline, in which LLM workers
generate statements and proof scripts and the Lean 4 kernel is the sole
adjudicator at every gate: statements are frozen by hash before proving, and a
separate audit re-computes the symmetric difference, scans the whole file for
\texttt{sorry}/\texttt{admit} including comments, and re-prints the axiom
dependencies \cite{lean4}. Verification strength is stated below verbatim.
This note is a priority claim of record deposited so that a dated, citable
public record exists ahead of the narrative paper of the parent program (the
defective-prism matching slices); that paper is prepared separately and will
reference this deposit. No literature review is attempted here.

\section{Statement}\label{sec:statement}

Throughout, graphs are finite simple graphs on a decidable type of vertices;
\texttt{hosoya} $G$ denotes $|M(G)|$, the number of matchings of $G$, defined
as the cardinality of the set of edge-subsets whose elements pairwise have
disjoint endpoint sets (the empty matching included). The prism graph is
$\mathtt{prism}\ n=(C_n\square P_2)$ on vertices $\texttt{Fin}\ n\times
\texttt{Fin}\ 2$; \texttt{prismDelSpoke} $n$ is the prism with the spoke edge
$((0,0),(0,1))$ deleted, endpoints kept (requirement $3\le n$ so the deletion
is well typed).

% LEAN: tasks/20260929-l14c-graphside-pipeline :: hosoya_deleteEdge_split_pts grade=完全证明
\begin{lemma}[edge-deletion matching-split identity; general form]\label{lem:split}
Let $G$ be a finite simple graph with decidable adjacency, and let $a,b$ be
adjacent vertices. Then
\[
\mathtt{hosoya}\ G\;=\;\mathtt{hosoya}\,(G.\mathtt{deleteEdges}\,\{ab\})
\;+\;\mathtt{hosoya}\,\bigl(G.\mathtt{induce}\,(\text{all vertices except }
a,b)\bigr).
\]
\end{lemma}

\begin{lstlisting}[caption={Frozen base definitions, verbatim lines 45--62 of
\texttt{formalized/statement.lean} (sha256 \texttt{d186ff16}): edgeSupp,
hosoya, prism, prismDelSpoke and their decidability instances.},label={lst:s1}]
def edgeSupp {V : Type*} [DecidableEq V] [Fintype V] (e : Sym2 V) : Finset V :=
  (univ : Finset V).filter fun v => v ∈ e

def hosoya (G : SimpleGraph V) [Fintype V] [DecidableEq V] [DecidableRel G.Adj] : ℕ :=
  ((G.edgeFinset.powerset).filter fun s =>
    (∀ e ∈ s, ∀ f ∈ s, e ≠ f → Disjoint (edgeSupp e) (edgeSupp f))).card

def prism (n : ℕ) : SimpleGraph (Fin n × Fin 2) := cycleGraph n □ pathGraph 2

instance (n : ℕ) : DecidableRel (prism n).Adj :=
  show DecidableRel (cycleGraph n □ pathGraph 2).Adj from inferInstance

def prismDelSpoke (n : ℕ) (hn : 3 ≤ n) : SimpleGraph (Fin n × Fin 2) :=
  (prism n).deleteEdges {Sym2.mk ((⟨0, by omega⟩, (0 : Fin 2))) ((⟨0, by omega⟩, (1 : Fin 2)))}

instance (n : ℕ) (hn : 3 ≤ n) : DecidableRel (prismDelSpoke n hn).Adj :=
  show DecidableRel ((prism n).deleteEdges
    {Sym2.mk ((⟨0, by omega⟩, (0 : Fin 2))) ((⟨0, by omega⟩, (1 : Fin 2)))}).Adj from inferInstance
\end{lstlisting}

\begin{lstlisting}[caption={General lemma, verbatim head of
\texttt{attempts/phase0/r2/split-identity.lean} lines 215--259 (declaration
and first half of the proof).},label={lst:s2}]
theorem hosoya_deleteEdge_split_pts {V : Type*} [Fintype V] [DecidableEq V]
    (G : SimpleGraph V) [DecidableRel G.Adj]
    (a b : V) (huv : G.Adj a b) :
    hosoya G = hosoya (G.deleteEdges {Sym2.mk a b}) +
      hosoya (SimpleGraph.induce (Set.univ \ {a, b} : Set V) G) := by
  have he : Sym2.mk a b ∈ G.edgeFinset := by
    rw [SimpleGraph.mem_edgeFinset, SimpleGraph.mem_edgeSet]; exact huv
  have hA : ((G.edgeFinset.powerset).filter fun s =>
        (∀ e₁ ∈ s, ∀ f ∈ s, e₁ ≠ f → Disjoint (edgeSupp e₁) (edgeSupp f)) ∧ Sym2.mk a b ∈ s).card
      = ((G.edgeFinset.powerset).filter fun u =>
        (∀ a₁ ∈ u, ∀ b₁ ∈ u, a₁ ≠ b₁ → Disjoint (edgeSupp a₁) (edgeSupp b₁)) ∧
          ∀ f ∈ u, Disjoint (edgeSupp f) (edgeSupp (Sym2.mk a b))).card := by
    apply Finset.card_bij (fun s _ => s.erase (Sym2.mk a b))
    · intro s hs
      rw [Finset.mem_filter, Finset.mem_powerset] at hs
      obtain ⟨hsE, hsP, hse⟩ := hs
      rw [Finset.mem_filter, Finset.mem_powerset]
      refine ⟨(Finset.erase_subset _ s).trans hsE, ?_, ?_⟩
      · intro a₁ ha b₁ hb hab
        exact hsP a₁ (Finset.mem_of_mem_erase ha) b₁ (Finset.mem_of_mem_erase hb) hab
      · intro f hf
        exact hsP f (Finset.mem_of_mem_erase hf) (Sym2.mk a b) hse (Finset.ne_of_mem_erase hf)
    · intro s₁ hs₁ s₂ hs₂ hh
      rw [Finset.mem_filter] at hs₁ hs₂
      rw [← Finset.insert_erase hs₁.2.2, ← Finset.insert_erase hs₂.2.2, hh]
    · intro u hu
      rw [Finset.mem_filter, Finset.mem_powerset] at hu
      obtain ⟨huE, huP, hdj⟩ := hu
      have heu : Sym2.mk a b ∉ u := by
        intro heu'
        have h := hdj (Sym2.mk a b) heu'
        obtain ⟨v, hv⟩ := Finset.card_pos.mp (by rw [edgeSupp_card_two G (Sym2.mk a b) he]; norm_num)
        exact Finset.disjoint_left.mp h hv hv
      refine ⟨insert (Sym2.mk a b) u, ?_, Finset.erase_insert heu⟩
      rw [Finset.mem_filter, Finset.mem_powerset]
      refine ⟨Finset.insert_subset he huE, ?_, Finset.mem_insert_self _ _⟩
      intro a₁ ha b₁ hb hab
      rw [Finset.mem_insert] at ha hb
      rcases ha with rfl | ha
      · rcases hb with rfl | hb
        · exact absurd rfl hab
        · exact (hdj b₁ hb).symm
      · rcases hb with rfl | hb
        · exact hdj a₁ ha
        · exact huP a₁ ha b₁ hb hab
\end{lstlisting}

% LEAN: tasks/20260929-l14c-graphside-pipeline :: hosoya_prism_split_spoke grade=完全证明
\begin{theorem}[instantiation on the prism graph]\label{thm:prismsplit}
For every $n$ with $3\le n$, the matching counts of the $n$-prism split along
the spoke edge $((0,0),(0,1))$:
\begin{align*}
\mathtt{hosoya}\,(\mathtt{prism}\ n)
={}&\mathtt{hosoya}\,(\mathtt{prismDelSpoke}\ n)\\
&+\mathtt{hosoya}\,(\text{prism induced on the two points' complement}).
\end{align*}
The induced factor is the path $P_{n-1}$ (numerically the ladder values
$8,21,...$); its identification with $P_{n-1}$ is not kernel-proved here and
is not claimed.
\end{theorem}

The complete final file is reproduced verbatim in
Section~\ref{sec:proof} (whole file, no excerpting); it also ships in this
deposit under \texttt{proofs/}. The frozen statement file
\texttt{formalized/statement.lean} of the underlying task also contains the
program's two main bridge theorems (defective-prism matchings equal the
recurrence-defined sequences) with \texttt{sorry} placeholders; they are
\emph{not} results of this deposit and are mentioned only to delimit its
scope. The statement excerpt above stops at the base definitions, so no
\texttt{sorry} line enters any listing.

\section{Proof}\label{sec:proof}

The proof of Lemma~\ref{lem:split} partitions the power set of the edge
finset by ``contains $ab$'' / ``avoids $ab$'' and counts each side: the
avoiding side is the matching filter on the edge finset with the edge removed
(a definitional rewriting of \texttt{deleteEdges}); the containing side is
mapped bijectively, via \texttt{Finset.card\_bij}, to the matching filter of
the induced graph on the two-point complement, with each $e$-free matching of
$G$ lifted edge by edge across the vertex-type change using \texttt{Sym2.map}
and its disjointness/injectivity transport lemmas. The instantiation
Theorem~\ref{thm:prismsplit} applies the general lemma to
$G=\mathtt{prism}\ n$ with $ab$ the spoke, discharging the member
adjacency through the box-product adjacency rule. The complete final file
follows verbatim ($429$ lines, no excerpting); it also ships in this deposit
under \texttt{proofs/}.

\begin{lstlisting}[basicstyle=\ttfamily\footnotesize,caption={Complete proof
file of Lemma~\ref{lem:split} and Theorem~\ref{thm:prismsplit}, verbatim lines
1--429 (whole file) of \texttt{attempts/phase0/r2/split-identity.lean} (sha256
\texttt{d2efac51}).},label={lst:proof}]
import Mathlib

open SimpleGraph Finset

/-!
# Phase 0 删边分割恒等式 · 完全证明（2026-09-29，零占位、严格编译通过）

 provenance：
 - 基座代码逐字承冻结快照 `formalized/statement.lean`（宪条 2，只许拷贝不许改）。
 - 「已验证块 A–D」逐字承 `attempts/phase0/r1/scratch-probe1.lean` 与
   `scratch-probe2.lean`（军政部 headless 工人 r1b 留下，主代理 `lake env lean`
   实测两文件均 exit 0 / 0 error，2026-09-29 15:2x）；本文件把两探针的 example
   转成可引用的 theorem（证明体逐字未动；`edgeSupp_card_two` 前加 `include G he`
   一处——Lean 4 section 变量只按声明签名收绑定件，证明体里单独用要 include）。
 - 一般装置引理 `hosoya_deleteEdge_split_pts`（两点集形态）：数学骨架承工人 r2 在
   `r2-scratch.lean` 中证明的 edgeSp 形态版（主代理实测编译通过），由主代理适配为
   两点集形态——原因：`Set.univ \ ↑(edgeSupp …)` 经 Finset coercion，TC 合成
   `Set.univ.fintypeDiffLeft` 实例，与变量集形态的通用 `Subtype.fintype` 实例
   非 defeq（实测 8M 心跳 type mismatch），实例化时 rw/exact 全线失配；
   两点集字面量与冻结陈述同一条实例合成路径（pts 探针实证 exit 0）。
 - 目标定理 `hosoya_prism_split_spoke`：statement 逐字冻结，实例化走两点集引理
   + `boxProd_adj` 邻接路线（承工人 P2 探针）。
-/

-- ===== 基座代码（逐字承冻结快照 formalized/statement.lean，只许拷贝不许改） =====

instance (n : ℕ) : DecidableRel (pathGraph n).Adj := fun _ _ =>
  decidable_of_iff _ pathGraph_adj.symm

instance {α β : Type*} [DecidableEq α] [DecidableEq β] (G : SimpleGraph α) (H : SimpleGraph β)
    [DecidableRel G.Adj] [DecidableRel H.Adj] : DecidableRel (G □ H).Adj := fun _ _ =>
  decidable_of_iff _ boxProd_adj.symm

def edgeSupp {V : Type*} [DecidableEq V] [Fintype V] (e : Sym2 V) : Finset V :=
  (univ : Finset V).filter fun v => v ∈ e

def hosoya (G : SimpleGraph V) [Fintype V] [DecidableEq V] [DecidableRel G.Adj] : ℕ :=
  ((G.edgeFinset.powerset).filter fun s =>
    (∀ e ∈ s, ∀ f ∈ s, e ≠ f → Disjoint (edgeSupp e) (edgeSupp f))).card

def prism (n : ℕ) : SimpleGraph (Fin n × Fin 2) := cycleGraph n □ pathGraph 2

instance (n : ℕ) : DecidableRel (prism n).Adj :=
  show DecidableRel (cycleGraph n □ pathGraph 2).Adj from inferInstance

def prismDelSpoke (n : ℕ) (hn : 3 ≤ n) : SimpleGraph (Fin n × Fin 2) :=
  (prism n).deleteEdges {Sym2.mk ((⟨0, by omega⟩, (0 : Fin 2))) ((⟨0, by omega⟩, (1 : Fin 2)))}

instance (n : ℕ) (hn : 3 ≤ n) : DecidableRel (prismDelSpoke n hn).Adj :=
  show DecidableRel ((prism n).deleteEdges
    {Sym2.mk ((⟨0, by omega⟩, (0 : Fin 2))) ((⟨0, by omega⟩, (1 : Fin 2)))}).Adj from inferInstance

-- ===== 已验证块 A（承 scratch-probe1.lean，rfl 展开层） =====

theorem hosoya_def {V : Type*} [DecidableEq V] [Fintype V] (G : SimpleGraph V)
    [DecidableRel G.Adj] :
    hosoya G = ((G.edgeFinset.powerset).filter fun s =>
      (∀ e ∈ s, ∀ f ∈ s, e ≠ f → Disjoint (edgeSupp e) (edgeSupp f))).card := rfl

-- ===== 已验证块 B（承 scratch-probe1.lean P2/P4，棱柱 spoke 事实） =====

theorem spoke_mem_edgeFinset (n : ℕ) (hn : 3 ≤ n) :
    Sym2.mk ((⟨0, by omega⟩ : Fin n), (0 : Fin 2))
      ((⟨0, by omega⟩ : Fin n), (1 : Fin 2)) ∈ (prism n).edgeFinset := by
  rw [SimpleGraph.mem_edgeFinset, SimpleGraph.mem_edgeSet]
  show (cycleGraph n □ pathGraph 2).Adj (⟨0, by omega⟩, (0 : Fin 2)) (⟨0, by omega⟩, (1 : Fin 2))
  rw [boxProd_adj]
  refine Or.inr ⟨?_, rfl⟩
  show (pathGraph 2).Adj (0 : Fin 2) (1 : Fin 2)
  decide

theorem edgeSupp_spoke (n : ℕ) (hn : 3 ≤ n) :
    (↑(edgeSupp (Sym2.mk ((⟨0, by omega⟩ : Fin n), (0 : Fin 2))
      ((⟨0, by omega⟩ : Fin n), (1 : Fin 2)))) : Set (Fin n × Fin 2)) =
    {((⟨0, by omega⟩ : Fin n), (0 : Fin 2)), ((⟨0, by omega⟩ : Fin n), (1 : Fin 2))} := by
  ext v
  simp [edgeSupp, Sym2.mem_iff]

-- ===== 已验证块 C（承 scratch-probe1.lean P5/P6，不含 e 半 + 分区计数） =====

theorem matching_filter_not_mem (V : Type*) [Fintype V] [DecidableEq V] (G : SimpleGraph V)
    [DecidableRel G.Adj] (e : Sym2 V) :
    (G.edgeFinset.powerset.filter fun s =>
        (∀ e₁ ∈ s, ∀ f ∈ s, e₁ ≠ f → Disjoint (edgeSupp e₁) (edgeSupp f)) ∧ e ∉ s) =
      ((G.edgeFinset \ {e}).powerset).filter fun s =>
        (∀ e₁ ∈ s, ∀ f ∈ s, e₁ ≠ f → Disjoint (edgeSupp e₁) (edgeSupp f)) := by
  ext s
  simp only [Finset.mem_powerset, Finset.mem_filter]
  constructor
  · rintro ⟨hs, hP, hne⟩
    refine ⟨fun x hx => ?_, hP⟩
    rw [Finset.mem_sdiff, Finset.mem_singleton]
    exact ⟨hs hx, fun h => hne (h ▸ hx)⟩
  · rintro ⟨hs, hP⟩
    refine ⟨fun x hx => (Finset.mem_sdiff.mp (hs hx)).1, hP, fun h => ?_⟩
    exact (Finset.mem_sdiff.mp (hs h)).2 (Finset.mem_singleton_self e)

theorem hosoya_partition (V : Type*) [Fintype V] [DecidableEq V] (G : SimpleGraph V)
    [DecidableRel G.Adj] (e : Sym2 V) :
    hosoya G = ((G.edgeFinset.powerset.filter fun s =>
        (∀ e₁ ∈ s, ∀ f ∈ s, e₁ ≠ f → Disjoint (edgeSupp e₁) (edgeSupp f)) ∧ e ∉ s).card +
      (G.edgeFinset.powerset.filter fun s =>
        (∀ e₁ ∈ s, ∀ f ∈ s, e₁ ≠ f → Disjoint (edgeSupp e₁) (edgeSupp f)) ∧ e ∈ s).card) := by
  have h1 := Finset.card_filter_add_card_filter_not
    (s := G.edgeFinset.powerset.filter fun s =>
      (∀ e₁ ∈ s, ∀ f ∈ s, e₁ ≠ f → Disjoint (edgeSupp e₁) (edgeSupp f))) (fun s => e ∉ s)
  rw [hosoya_def, ← h1]
  congr 1
  · rw [Finset.filter_filter]
  · rw [Finset.filter_filter]
    congr 1
    ext s
    simp only [not_not]

-- ===== 已验证块 D（承 scratch-probe2.lean，Sym2 跨顶点型搬运四件） =====

theorem edgeSupp_mem {V : Type*} [DecidableEq V] [Fintype V] (e : Sym2 V) (v : V) :
    v ∈ edgeSupp e ↔ v ∈ e := by simp [edgeSupp]

section Transport

variable {V : Type*} [Fintype V] [DecidableEq V] (G : SimpleGraph V) [DecidableRel G.Adj]
  (e : Sym2 V) (he : e ∈ G.edgeFinset)

theorem edgeSupp_map_subtype (g : Sym2 ↥(↑((univ : Finset V) \ edgeSupp e) : Set V)) :
    (edgeSupp g).map (Function.Embedding.subtype
      (· ∈ (↑((univ : Finset V) \ edgeSupp e) : Set V))) =
    edgeSupp (Sym2.map (Subtype.val) g) := by
  ext v
  rw [Finset.mem_map]
  constructor
  · rintro ⟨u, hu, huv⟩
    rw [edgeSupp_mem] at hu
    rw [edgeSupp_mem]
    exact Sym2.mem_map.mpr ⟨u, hu, huv⟩
  · intro hv
    rw [edgeSupp_mem] at hv
    obtain ⟨u, hu, huv⟩ := Sym2.mem_map.mp hv
    exact ⟨u, by rw [edgeSupp_mem]; exact hu, huv⟩

theorem disjoint_map_subtype (g₁ g₂ : Sym2 ↥(↑((univ : Finset V) \ edgeSupp e) : Set V)) :
    Disjoint (edgeSupp g₁) (edgeSupp g₂) ↔
    Disjoint (edgeSupp (Sym2.map (Subtype.val) g₁)) (edgeSupp (Sym2.map (Subtype.val) g₂)) := by
  have L1 : ∀ g : Sym2 ↥(↑((univ : Finset V) \ edgeSupp e) : Set V),
      (edgeSupp g).map (Function.Embedding.subtype
        (· ∈ (↑((univ : Finset V) \ edgeSupp e) : Set V))) =
      edgeSupp (Sym2.map (Subtype.val) g) := by
    intro g
    ext v
    rw [Finset.mem_map]
    constructor
    · rintro ⟨u, hu, huv⟩
      rw [edgeSupp_mem] at hu
      rw [edgeSupp_mem]
      exact Sym2.mem_map.mpr ⟨u, hu, huv⟩
    · intro hv
      rw [edgeSupp_mem] at hv
      obtain ⟨u, hu, huv⟩ := Sym2.mem_map.mp hv
      exact ⟨u, by rw [edgeSupp_mem]; exact hu, huv⟩
  rw [← L1 g₁, ← L1 g₂]
  constructor
  · intro h
    rw [Finset.disjoint_left] at h ⊢
    intro v hv hv'
    rw [Finset.mem_map] at hv hv'
    obtain ⟨u, hu, huv⟩ := hv
    obtain ⟨u', hu', huv'⟩ := hv'
    have huu : u = u' := (Function.Embedding.subtype _).injective (huv.trans huv'.symm)
    exact h hu (huu ▸ hu')
  · intro h
    rw [Finset.disjoint_left] at h ⊢
    intro u hu hu'
    exact h (Finset.mem_map.mpr ⟨u, hu, rfl⟩) (Finset.mem_map.mpr ⟨u, hu', rfl⟩)

theorem mapped_ne_edge (g : Sym2 ↥(↑((univ : Finset V) \ edgeSupp e) : Set V)) :
    Sym2.map (Subtype.val) g ≠ e := by
  intro hge
  have hex : ∃ u : V, u ∈ e := by
    induction e using Sym2.ind with
    | _ a b => exact ⟨a, Sym2.mem_iff.mpr (Or.inl rfl)⟩
  obtain ⟨u, hu⟩ := hex
  have huE : u ∈ edgeSupp e := Finset.mem_filter.mpr ⟨Finset.mem_univ u, hu⟩
  have hu2 : u ∈ edgeSupp (Sym2.map (Subtype.val) g) := hge.symm ▸ huE
  rw [edgeSupp_mem] at hu2
  obtain ⟨w, hwg, hwu⟩ := Sym2.mem_map.mp hu2
  rw [← hwu] at huE
  have hwT : (↑w : V) ∈ (↑((univ : Finset V) \ edgeSupp e) : Set V) := w.2
  rw [Finset.mem_coe, Finset.mem_sdiff] at hwT
  exact hwT.2 huE

include G he
theorem edgeSupp_card_two : (edgeSupp e).card = 2 := by
  rw [SimpleGraph.mem_edgeFinset] at he
  induction e using Sym2.ind with
  | _ a b =>
    rw [SimpleGraph.mem_edgeSet] at he
    have hab : a ≠ b := G.ne_of_adj he
    rw [Finset.card_eq_two]
    refine ⟨a, b, hab, ?_⟩
    ext v
    simp [edgeSupp, Sym2.mem_iff]

end Transport

-- ===== 一般装置引理（删边分割恒等式 · 两点集形态） =====

/-- 一般装置引理（删边分割恒等式，两点集形态）：对任意图 `G` 与边 `a ~ b`，
匹配集按「含 `s(a, b)` / 不含 `s(a, b)`」二分——不含者 = `G.deleteEdges {s(a,b)}`
的匹配数（已验证块 C），含者 = 删 `a, b` 两端点后诱导子图的匹配数（B 侧：
`Finset.card_bij` 双豫 + Sym2 跨型搬运）。
形态说明：诱导集用 `Set.univ \ {a, b}`（两点集字面量）而非 `Set.univ \ ↑(edgeSupp …)`——
后者经 Finset coercion，TC 合成的是 `Set.univ.fintypeDiffLeft` 实例，与变量集形态的
通用 `Subtype.fintype` 实例非 defeq（实测 8M 心跳 type mismatch），会导致实例化时
rw/exact 全线失配；两点集字面量形态与冻结陈述同一条实例合成路径（pts 探针实证）。 -/
theorem hosoya_deleteEdge_split_pts {V : Type*} [Fintype V] [DecidableEq V]
    (G : SimpleGraph V) [DecidableRel G.Adj]
    (a b : V) (huv : G.Adj a b) :
    hosoya G = hosoya (G.deleteEdges {Sym2.mk a b}) +
      hosoya (SimpleGraph.induce (Set.univ \ {a, b} : Set V) G) := by
  have he : Sym2.mk a b ∈ G.edgeFinset := by
    rw [SimpleGraph.mem_edgeFinset, SimpleGraph.mem_edgeSet]; exact huv
  have hA : ((G.edgeFinset.powerset).filter fun s =>
        (∀ e₁ ∈ s, ∀ f ∈ s, e₁ ≠ f → Disjoint (edgeSupp e₁) (edgeSupp f)) ∧ Sym2.mk a b ∈ s).card
      = ((G.edgeFinset.powerset).filter fun u =>
        (∀ a₁ ∈ u, ∀ b₁ ∈ u, a₁ ≠ b₁ → Disjoint (edgeSupp a₁) (edgeSupp b₁)) ∧
          ∀ f ∈ u, Disjoint (edgeSupp f) (edgeSupp (Sym2.mk a b))).card := by
    apply Finset.card_bij (fun s _ => s.erase (Sym2.mk a b))
    · intro s hs
      rw [Finset.mem_filter, Finset.mem_powerset] at hs
      obtain ⟨hsE, hsP, hse⟩ := hs
      rw [Finset.mem_filter, Finset.mem_powerset]
      refine ⟨(Finset.erase_subset _ s).trans hsE, ?_, ?_⟩
      · intro a₁ ha b₁ hb hab
        exact hsP a₁ (Finset.mem_of_mem_erase ha) b₁ (Finset.mem_of_mem_erase hb) hab
      · intro f hf
        exact hsP f (Finset.mem_of_mem_erase hf) (Sym2.mk a b) hse (Finset.ne_of_mem_erase hf)
    · intro s₁ hs₁ s₂ hs₂ hh
      rw [Finset.mem_filter] at hs₁ hs₂
      rw [← Finset.insert_erase hs₁.2.2, ← Finset.insert_erase hs₂.2.2, hh]
    · intro u hu
      rw [Finset.mem_filter, Finset.mem_powerset] at hu
      obtain ⟨huE, huP, hdj⟩ := hu
      have heu : Sym2.mk a b ∉ u := by
        intro heu'
        have h := hdj (Sym2.mk a b) heu'
        obtain ⟨v, hv⟩ := Finset.card_pos.mp (by rw [edgeSupp_card_two G (Sym2.mk a b) he]; norm_num)
        exact Finset.disjoint_left.mp h hv hv
      refine ⟨insert (Sym2.mk a b) u, ?_, Finset.erase_insert heu⟩
      rw [Finset.mem_filter, Finset.mem_powerset]
      refine ⟨Finset.insert_subset he huE, ?_, Finset.mem_insert_self _ _⟩
      intro a₁ ha b₁ hb hab
      rw [Finset.mem_insert] at ha hb
      rcases ha with rfl | ha
      · rcases hb with rfl | hb
        · exact absurd rfl hab
        · exact (hdj b₁ hb).symm
      · rcases hb with rfl | hb
        · exact hdj a₁ ha
        · exact huP a₁ ha b₁ hb hab
  have hB : ((G.edgeFinset.powerset).filter fun u =>
        (∀ a₁ ∈ u, ∀ b₁ ∈ u, a₁ ≠ b₁ → Disjoint (edgeSupp a₁) (edgeSupp b₁)) ∧
          ∀ f ∈ u, Disjoint (edgeSupp f) (edgeSupp (Sym2.mk a b))).card
      = (((SimpleGraph.induce (Set.univ \ {a, b} : Set V) G).edgeFinset.powerset).filter
          fun t => (∀ a₁ ∈ t, ∀ b₁ ∈ t, a₁ ≠ b₁ → Disjoint (edgeSupp a₁) (edgeSupp b₁))).card := by
    have hTS : ∀ v : V, v ∉ edgeSupp (Sym2.mk a b) → v ∈ (Set.univ \ {a, b} : Set V) := by
      intro v hv
      refine ⟨Set.mem_univ v, ?_⟩
      intro hmem
      apply hv
      rw [edgeSupp_mem, Sym2.mem_iff]
      simpa using hmem
    have map_inj : Function.Injective
        (Sym2.map (Subtype.val : ↥(Set.univ \ {a, b} : Set V) → V)) := by
      intro x y hxy
      induction x using Sym2.ind with
      | _ x₁ y₁ =>
        induction y using Sym2.ind with
        | _ x₂ y₂ =>
          rw [Sym2.map_mk, Sym2.map_mk, Sym2.eq_iff] at hxy
          rw [Sym2.eq_iff]
          rcases hxy with ⟨h1, h2⟩ | ⟨h1, h2⟩
          · exact Or.inl ⟨Subtype.ext h1, Subtype.ext h2⟩
          · exact Or.inr ⟨Subtype.ext h1, Subtype.ext h2⟩
    have hlift : ∀ f : Sym2 V, f ∈ G.edgeFinset → Disjoint (edgeSupp f) (edgeSupp (Sym2.mk a b)) →
        ∃ g : Sym2 ↥(Set.univ \ {a, b} : Set V),
          Sym2.map Subtype.val g = f ∧
          g ∈ (SimpleGraph.induce (Set.univ \ {a, b} : Set V) G).edgeFinset := by
      intro f hfE hfdj
      induction f using Sym2.ind with
      | _ x y =>
        rw [SimpleGraph.mem_edgeFinset, SimpleGraph.mem_edgeSet] at hfE
        have hxE : x ∉ edgeSupp (Sym2.mk a b) := by
          intro ha
          exact Finset.disjoint_left.mp hfdj
            (by rw [edgeSupp_mem]; exact Sym2.mem_iff.mpr (Or.inl rfl)) ha
        have hyE : y ∉ edgeSupp (Sym2.mk a b) := by
          intro hb
          exact Finset.disjoint_left.mp hfdj
            (by rw [edgeSupp_mem]; exact Sym2.mem_iff.mpr (Or.inr rfl)) hb
        refine ⟨s((⟨x, hTS x hxE⟩ : ↥(Set.univ \ {a, b} : Set V)),
          (⟨y, hTS y hyE⟩ : ↥(Set.univ \ {a, b} : Set V))), ?_, ?_⟩
        · rw [Sym2.map_mk]
        · rw [SimpleGraph.mem_edgeFinset, SimpleGraph.mem_edgeSet]
          exact SimpleGraph.induce_adj.mpr hfE
    symm
    apply Finset.card_bij (fun t _ => t.image (Sym2.map (Subtype.val :
      ↥(Set.univ \ {a, b} : Set V) → V)))
    · intro t ht
      rw [Finset.mem_filter, Finset.mem_powerset] at ht
      obtain ⟨htE, htP⟩ := ht
      rw [Finset.mem_filter, Finset.mem_powerset]
      refine ⟨?_, ?_, ?_⟩
      · intro y hy
        obtain ⟨g, hg, rfl⟩ := Finset.mem_image.mp hy
        have hgE := htE hg
        induction g using Sym2.ind with
        | _ x y =>
          rw [SimpleGraph.mem_edgeFinset, SimpleGraph.mem_edgeSet] at hgE
          rw [Sym2.map_mk, SimpleGraph.mem_edgeFinset, SimpleGraph.mem_edgeSet]
          exact SimpleGraph.induce_adj.mp hgE
      · intro x hx y₁ hy hxy
        obtain ⟨g₁, hg₁, rfl⟩ := Finset.mem_image.mp hx
        obtain ⟨g₂, hg₂, rfl⟩ := Finset.mem_image.mp hy
        have hgg : g₁ ≠ g₂ := fun h => hxy (by rw [h])
        have hdj := htP g₁ hg₁ g₂ hg₂ hgg
        rw [Finset.disjoint_left] at hdj ⊢
        intro v hv1 hv2
        rw [edgeSupp_mem] at hv1 hv2
        obtain ⟨u₁, hu₁, huv₁⟩ := Sym2.mem_map.mp hv1
        obtain ⟨u₂, hu₂, huv₂⟩ := Sym2.mem_map.mp hv2
        have huu : u₁ = u₂ := Subtype.ext (huv₁.trans huv₂.symm)
        exact hdj ((edgeSupp_mem g₁ u₁).mpr hu₁) ((edgeSupp_mem g₂ u₁).mpr (huu ▸ hu₂))
      · intro y hy
        obtain ⟨g, hg, rfl⟩ := Finset.mem_image.mp hy
        rw [Finset.disjoint_left]
        intro v hv hve
        rw [edgeSupp_mem] at hv
        obtain ⟨u, hu, huv⟩ := Sym2.mem_map.mp hv
        rw [← huv] at hve
        have huS := (Set.mem_sdiff _).mp u.2
        have hmem : (u.1 : V) ∈ ({a, b} : Set V) := by
          have h1 : u.1 ∈ Sym2.mk a b := (edgeSupp_mem _ _).mp hve
          have h2 : u.1 = a ∨ u.1 = b := (Sym2.mem_iff).mp h1
          simpa using h2
        exact huS.2 hmem
    · intro t₁ h₁ t₂ h₂ hh
      ext g
      constructor
      · intro hg
        have hm : Sym2.map Subtype.val g ∈ t₂.image (Sym2.map Subtype.val) := by
          rw [← hh]
          exact Finset.mem_image.mpr ⟨g, hg, rfl⟩
        obtain ⟨g', hg', hgg'⟩ := Finset.mem_image.mp hm
        exact (map_inj hgg') ▸ hg'
      · intro hg
        have hm : Sym2.map Subtype.val g ∈ t₁.image (Sym2.map Subtype.val) := by
          rw [hh]
          exact Finset.mem_image.mpr ⟨g, hg, rfl⟩
        obtain ⟨g', hg', hgg'⟩ := Finset.mem_image.mp hm
        exact (map_inj hgg') ▸ hg'
    · intro u hu
      rw [Finset.mem_filter, Finset.mem_powerset] at hu
      obtain ⟨huE, huP, hdj⟩ := hu
      refine ⟨u.attach.image fun x => Classical.choose (hlift x.1 (huE x.2) (hdj x.1 x.2)), ?_, ?_⟩
      · rw [Finset.mem_filter, Finset.mem_powerset]
        constructor
        · intro g hg
          obtain ⟨⟨f, hf⟩, -, rfl⟩ := Finset.mem_image.mp hg
          exact (Classical.choose_spec (hlift f (huE hf) (hdj f hf))).2
        · intro g₁ hg₁ g₂ hg₂ hgg
          obtain ⟨⟨f₁, hf₁⟩, -, rfl⟩ := Finset.mem_image.mp hg₁
          obtain ⟨⟨f₂, hf₂⟩, -, rfl⟩ := Finset.mem_image.mp hg₂
          have hff : f₁ ≠ f₂ := by
            intro h
            apply hgg
            subst h
            rfl
          have s1 := (Classical.choose_spec (hlift f₁ (huE hf₁) (hdj f₁ hf₁))).1
          have s2 := (Classical.choose_spec (hlift f₂ (huE hf₂) (hdj f₂ hf₂))).1
          have hdj2 := huP f₁ hf₁ f₂ hf₂ hff
          rw [Finset.disjoint_left] at hdj2 ⊢
          intro v hv1 hv2
          rw [edgeSupp_mem] at hv1 hv2
          exact hdj2
            ((edgeSupp_mem f₁ v).mpr (s1 ▸ Sym2.mem_map.mpr ⟨v, hv1, rfl⟩))
            ((edgeSupp_mem f₂ v).mpr (s2 ▸ Sym2.mem_map.mpr ⟨v, hv2, rfl⟩))
      · ext z
        constructor
        · intro hz
          obtain ⟨g, hg, hgz⟩ := Finset.mem_image.mp hz
          obtain ⟨⟨f, hf⟩, -, rfl⟩ := Finset.mem_image.mp hg
          have s := (Classical.choose_spec (hlift f (huE hf) (hdj f hf))).1
          rw [s] at hgz
          exact hgz ▸ hf
        · intro hz
          have s := (Classical.choose_spec (hlift z (huE hz) (hdj z hz))).1
          exact Finset.mem_image.mpr ⟨_, Finset.mem_image.mpr ⟨⟨z, hz⟩, Finset.mem_attach _ _, rfl⟩, s⟩
  have h2 : (G.deleteEdges {Sym2.mk a b}).edgeFinset = G.edgeFinset \ {Sym2.mk a b} := by
    ext f
    induction f using Sym2.ind with
    | _ x y =>
      simp only [Finset.mem_sdiff, Finset.mem_singleton]
      rw [SimpleGraph.mem_edgeFinset, SimpleGraph.mem_edgeSet,
        SimpleGraph.mem_edgeFinset, SimpleGraph.mem_edgeSet, SimpleGraph.deleteEdges_adj,
        Set.mem_singleton_iff]
  have hdel : hosoya (G.deleteEdges {Sym2.mk a b}) =
      ((G.edgeFinset \ {Sym2.mk a b}).powerset.filter fun s =>
        (∀ e₁ ∈ s, ∀ f ∈ s, e₁ ≠ f → Disjoint (edgeSupp e₁) (edgeSupp f))).card := by
    rw [hosoya_def, h2]
  rw [hosoya_partition V G (Sym2.mk a b), matching_filter_not_mem V G (Sym2.mk a b), hdel,
    hosoya_def, hA, hB]
  congr 1

-- ===== 目标定理（statement 逐字冻结；实例化走两点集形态引理） =====

theorem hosoya_prism_split_spoke : ∀ (n : ℕ) (hn : 3 ≤ n),
    hosoya (prism n) =
      hosoya (prismDelSpoke n hn) +
        hosoya (SimpleGraph.induce
          (Set.univ \ {((⟨0, by omega⟩, (0 : Fin 2))), ((⟨0, by omega⟩, (1 : Fin 2)))} : Set (Fin n × Fin 2))
          (prism n)) := by
  intro n hn
  refine hosoya_deleteEdge_split_pts (prism n) ((⟨0, by omega⟩, (0 : Fin 2)))
    ((⟨0, by omega⟩, (1 : Fin 2))) ?_
  show (cycleGraph n □ pathGraph 2).Adj (⟨0, by omega⟩, (0 : Fin 2)) (⟨0, by omega⟩, (1 : Fin 2))
  rw [boxProd_adj]
  refine Or.inr ⟨?_, rfl⟩
  show (pathGraph 2).Adj (0 : Fin 2) (1 : Fin 2)
  decide
\end{lstlisting}

\section{Verification evidence}

\begin{itemize}
\item Toolchain: Lean \texttt{v4.34.0} (\texttt{leanprover/lean4:v4.34.0}),
mathlib4 \texttt{v4.34.0}, revision \texttt{5ed29652}. Reproduction: with that
pin, \texttt{lake env lean} on \texttt{proofs/01-hosoya-split-proved.lean}
compiles with exit code $0$ and empty output.
\item Statement integrity: the frozen statement file
\texttt{formalized/statement.lean} ($124$ lines) has sha256
\begin{center}
\texttt{d186ff1651a4383e7e213b42ba4ec8c01}\allowbreak
\texttt{71d9e7004231db2764370cd967ca39e},
\end{center}
and the final file \texttt{split-identity.lean} ($429$ lines) has sha256
\begin{center}
\texttt{d2efac514083c505b9fa14675e59a8009}\allowbreak
\texttt{054e2063769c263e069d608f656b016}.
\end{center}
The theorem head of \texttt{hosoya\_prism\_split\_spoke} is byte-identical to
the frozen statement (programmatic comparison, zero symmetric difference; see
\texttt{audit/statement-diff.out}). The base definitions excerpted in
Listing~\ref{lst:s1} are copied verbatim from the frozen statement file.
\item Kernel check: compiles with $0$ \texttt{sorry} (whole-file scan
including comments, case-insensitive, also covering \texttt{admit},
\texttt{sorryAx}, \texttt{native\_decide}, \texttt{axiom} and
\texttt{unsafe}: zero hits); \texttt{\#print axioms} output, verbatim (kept
literal; the smaller monospace width is typesetting only):
\begin{lstlisting}[basicstyle=\ttfamily\footnotesize,frame=none]
'hosoya_deleteEdge_split_pts' depends on axioms: [propext, Classical.choice, Quot.sound]
'hosoya_prism_split_spoke' depends on axioms: [propext, Classical.choice, Quot.sound]
\end{lstlisting}
Each line is a subset of the whitelist \{propext, Quot.sound,
Classical.choice\}. The strict re-compile emits no error and no warning
(stdout empty); the evidence ships as \texttt{audit/audit-strict-compile.log}
and \texttt{audit/audit-axioms-sidecar.log}.
\item Grading by the audit department: both proven declarations
\emph{complete proof}; lowest grade reached for the case = \emph{complete
proof}, with one provenance caveat stated by the audit and carried here: the
mathematical skeleton of Lemma~\ref{lem:split} was produced by a worker agent
in a scratch probe (compiling) and was adapted to the two-point-set form by
the main agent to match the frozen statement's instance-synthesis path; the
provenance is recorded in the final file header and in the audit record, and
the deposit does not narrate the lemma as an independent single-author pass.
\item Audit chain (originals in this deposit under \texttt{audit/} unless
noted): final adjudication \texttt{final-01-hosoya-split.txt} (four checks:
statement comparison, strict compile, sorry scan, axiom audit), the statement
symmetric-difference record \texttt{statement-diff.out}, and this note's
claim check \texttt{claim-check.md} --- added to \texttt{audit/} after the
audit department returns its adjudication, and a prerequisite for
publication. The semantic-fidelity check (roundtrip re-translation) of the
underlying task ran on the same model node as its formalisation; that
independence weakness is mitigated, not removed, by kernel adjudication.
\end{itemize}

\section{Novelty evidence}\label{sec:novelty}

Adjudication copied from the pipeline's exclusivity review (C9 gate), which
ran before the problem was accepted and re-derived the data independently:
verdict \textbf{no occupying record found (reachable-channel scope)} for both
defective-prism matching slices and for the graph-side statement layer; value
tier of the parent topic \textbf{new sequence} (the mother sequence
A102080 \cite{a102080} is occupied and is not claimed), with the graph-side
bridge graded \textbf{first formalisation (Lean three-layer zero same-form)}.
The identity itself is not new mathematics and is not asserted as such: as a
human-maths tool it is classical in the matching-polynomial literature, and
the claim is strictly \emph{no prior Lean formalisation was reachable by the
review}. Summary by channel; the full query log, including every zero-hit
query, is in this deposit at \texttt{audit/c9-record.md}.

\begin{center}
\resizebox{\textwidth}{!}{%
\begin{tabular}{lll}
\hline
channel & queries & result \\
\hline
OEIS (dossier, sequence side) \cite{oeis} & 14 strings, 4 positive controls &
two slice families in six forms: zero hits; controls hit A102080 only \\
OEIS (graph-side increment) & signature $(2,4,0,-1)$ index-page sweep &
4 entries registered, neither slice among them; mother A102080 confirmed \\
arXiv & 9 (dossier + increment) & nearest neighbours read full-text
(PH-property in graph prisms; maritime keyword collision): no matching-count
or edge-deletion same-form \\
arXiv (packager's supplementary sweep, 2026-09-29) & all:"Hosoya index" AND
all:prism; abs:"Hosoya index" AND abs:"deleted"/"edge deletion" & 0 entries \\
Crossref & 7 & nearest hit Aldred--Plummer (matching extension):
structural, not counting; full text unread, recorded \\
MSE & 5 + 1 control & control hit; slice queries zero \\
GitHub code search (public Lean ecosystem) & 6 & zero same-form;
prism/hosoya/matching-count intersections all empty \\
mathlib rg layer (pinned v4.34.0) & 5 & matching predicates only, no
counting API; no prism definition; no Hosoya word family \\
independent review & 7 items & overall = confirmed (errata on two
non-load-bearing ledger blemishes) \\
\hline
\end{tabular}%
}
\end{center}

Two points of the record are part of the evidence. First, the review's
anchors self-proved (the mother string hitting A102080 with the exact name
line fixes what ``zero hit'' means on the sequence channel). Second, the
channels are named so a later reader can re-run them; if any channel later
returns a same-form Lean formalisation of the edge-deletion splitting
identity, the tier above weakens accordingly, and this deposit will be
amended by a later version rather than by editing this record.

\section{Scope and limitations}

Grade and boundary, stated as the audit states them. (a) The deposit proves
Lemma~\ref{lem:split} (general form) and Theorem~\ref{thm:prismsplit}
(prism instantiation) and nothing else. (b) The parent program's two bridge
theorems (defective-prism matchings equal recurrence-defined sequences) and
the six congruence-pattern theorems of the sequence side are \emph{not}
results of this deposit; the bridge theorems are stated with \texttt{sorry} in
the frozen statement file and their proof phases are in progress. (c)
Identification of the induced factor of Theorem~\ref{thm:prismsplit} with the
path $P_{n-1}$ is numerically supported but not kernel-proved here.
(d) \texttt{hosoya} is defined on decidable graphs; the identity is
finite-combinatorial and holds with or without decidability, but the lean
formulation fixes the decidable case. (e) The novelty record was scoped to
the defective-prism slices; the packager's supplementary arXiv sweeps of
2026-09-29 (Section~\ref{sec:novelty}) extend the zero-reachability evidence
to the identity's own wording, but channels outside the reachable set
(paywalled full texts, e.g.\ Farrell's ladder papers, and general web search)
remain recorded as unverified rather than as evidence. (f) The gate-four
sign-off of the underlying task was exercised by the author's instruction to
deposit this claim on 2026-09-29, while the task report file and the
department audit record are added to the working repository alongside this
deposit; the task's Phase 1/2 remain in flight and are not part of this
deposit. (g) This note carries no literature review and is a deposit record,
not a survey; the narrative paper of the parent program is separate.

\section*{Data availability}

The deposit package for this claim contains the claim note (LaTeX source and
PDF), the Lean sources (\texttt{proofs/}: frozen statement file, final proof
file, \texttt{lean-toolchain}), the verification and audit logs, and the full
C9 query log with zero-hit entries (\texttt{audit/}). Zenodo: the DOI is
reserved at upload by the author and is recorded in
\texttt{metadata/README.md} of this package once created; see
\texttt{UPLOAD.md} in the source repository. Claim note under CC BY 4.0, code
under MIT.

\section*{Use of generative AI}

(a) AI performed: candidate generation and the exclusivity review,
autoformalisation of the definitions and statements, the roundtrip semantic
check, proof search and proof-script writing (worker agents plus main-agent
form adaptation on Lemma~\ref{lem:split} as recorded above), error
classification and repair, and the assembly of this note including its
verbatim listings. (b) Model, node and budget: the proof phase ran on one
node, wire-id \texttt{anthropic/stepfun/step-5-preview}, with no cross-node
switching, consuming $2$ agent-run samples of a permitted $6$ for the phase
and $3$ kernel compiles of a permitted $36$ for the task, plus $8$
llm-call requests of a permitted $8$ ($4$ successful roundtrip adjudications
and $4$ transmission failures), copied verbatim from the task ledger. This
note consumed $0$ pipeline LLM calls against a cap of $4$; drafting was
mechanical assembly. (c) Final adjudication: all statements and proofs
reported here were checked by the Lean 4 kernel: the proof compiles with no
\texttt{sorry}, and \texttt{\#print axioms} reports only \{propext,
Classical.choice, Quot.sound\}. (d) The AI system is not an author; the human
author reviewed the evidence and is responsible for all content.

\section*{Author contributions}

CRediT-style: the human author adjudicated topic selection, statement
semantics and the deposit decision; the AI4Math pipeline (an LLM) generated
the candidate propositions, the proof searches and the text drafts. The AI
system is not listed as an author.

\begin{thebibliography}{9}
% Each entry verified on 2026-09-29 against Crossref title/page metadata and
% live fetches of oeis.org (A102080) and the arXiv export API. Raw responses:
% claims/hosoya-split-identity/audit/. No entry is quoted from memory.
\bibitem{lean4}
L.~de Moura and S.~Ullrich.
\newblock The Lean 4 theorem prover and programming language.
\newblock \emph{Automated Deduction -- CADE-28}, LNCS, pp.~625--635, 2021.
\url{https://doi.org/10.1007/978-3-030-79876-5_37}

\bibitem{mathlib}
The mathlib Community.
\newblock The Lean mathematical library.
\newblock \emph{CPP 2020}, pp.~367--381, 2020.
\url{https://doi.org/10.1145/3372885.3373824}

\bibitem{godsil-gutman}
C.~D. Godsil and I.~Gutman.
\newblock On the theory of the matching polynomial.
\newblock \emph{Journal of Graph Theory} 5 (1981), pp.~137--144.
\url{https://doi.org/10.1002/jgt.3190050203}

\bibitem{oeis}
The OEIS Foundation.
\newblock \emph{The On-Line Encyclopedia of Integer Sequences}.
\newblock \url{https://oeis.org/}

\bibitem{a102080}
OEIS Foundation.
\newblock Number of matchings in the $C_n\times P_2$ ($n$-prism) graph.
\newblock \emph{The On-Line Encyclopedia of Integer Sequences}, entry
A102080.
\url{https://oeis.org/A102080}
\end{thebibliography}

\end{document}